\section{Discussion}

\label{sec:discussion}
% ══════════════════════════════════════════════════════════════════════════════

\subsection{Comparison with Carbon Markets}

BIBs differ from carbon credits in key ways: BIBs are financial
instruments with contractual cash flows while credits are commodity
units; BIBs fund specific projects while credits are fungible; BIBs
have defined maturities and milestones while credits are static
certificates; and BIBs require continuous monitoring while credits use
periodic assessment.

\subsection{Scalability}

The BIB framework scales through template contracts reducing
per-issuance costs, shared oracle infrastructure serving multiple BIBs,
portfolio diversification reducing idiosyncratic risk, and index products
enabling passive conservation investment.

\subsection{Limitations}

\begin{enumerate}
  \item \textbf{Outcome payer risk}: Government commitment may be
    unreliable over 10+ year horizons.
  \item \textbf{Metric gaming}: Metric selection can inadvertently
    incentivize narrow optimization at expense of broader goals.
  \item \textbf{Technology risk}: Smart contract vulnerabilities, oracle
    manipulation, and blockchain scalability remain concerns.
  \item \textbf{Regulatory uncertainty}: Securities classification
    varies across jurisdictions.
  \item \textbf{Community displacement}: Outcome-based payments may
    incentivize exclusionary conservation; safeguards essential.
\end{enumerate}


% ══════════════════════════════════════════════════════════════════════════════
